\begin{afterword}
\summaryhead

The author once said that most organized belief systems exist to flee from doubt, and a listener objected that Jesuit novices are told to doubt their faith. The author could not confirm this, and treats it as a possibly hypothetical case. Such Jesuits would not be fleeing from doubt, the author concedes, but the practice looks suspicious: to a truly virtuous rationalist doubt should not be scary, and this sounds like getting used to something frightening.

The post then says what doubt is for. Like curiosity, a doubt aims at its own end. A rational doubt arises from a specific reason to suspect a belief, points to an investigation, and ends either by destroying the belief or by dying. A doubt that is not investigated or resolved does nothing. People admire doubt because they confuse it with modesty, and a display of doubt is worth no more than wearing a lab coat. The post ends with six rules, the last being that to an ideal mind doubt would not be scary.

\responsehead

The post's main point is sound: a doubt worth having comes from a specific reason and leads to an investigation, and doubt put on for show is only professed. It handles its opening anecdote fairly, stating that the report about the Jesuits could not be confirmed and treating the case as hypothetical throughout.

The main point is also old, and the post presents it as neglected. It says that ``it is not widely understood that you need a particular reason to doubt.'' This was Charles Peirce's argument against Descartes. In 1868 Peirce wrote that someone who comes to doubt a belief ``doubts because he has a positive reason for it,'' and ``Let us not pretend to doubt in philosophy what we do not doubt in our hearts.'' In 1877 he added that ``There must be a real and living doubt,'' and that inquiry begins with doubt and ends when doubt ends. Peirce is not mentioned.

The post's most absolute rule is overstated. ``An unresolved doubt is a null-op; it does not turn the wheel.'' The author's own ``The Proper Use of Humility,'' to which the post refers for a related point, praises an engineer who builds fail-safes into a machine while ``damn sure'' it will not fail. That engineer's doubt is not resolved, and it changes the design. Some doubts cannot be settled by investigation, and a reasonable response to them is a hedge. The rule fits the doubts that can be settled.

The judgment of the Jesuits sits uneasily with the post's own ending. The post finds a program of getting used to doubt ``highly suspicious,'' because doubt should not be scary. Its last rule grants that ``it may take courage to face your doubts,'' and says that only ``to an ideal mind'' doubt would not be scary. For minds that are not ideal, practice at facing frightening doubts is what that rule would seem to recommend.

The explanation of why people display doubt, that ``we all want to be seen as rational'' and confuse doubt with modesty, is given without evidence.

In short: a sound account of doubt, given by Peirce in 1868 and 1877 and here uncredited, with a rule about unresolved doubt that does not fit the author's own example of good humility.
\end{afterword}
